\documentclass[../../main.tex]{subfiles}
\begin{document}

\section{Route F: inverse-variance conditional-fiber frames}
\label{sec:conditional-fiber-frame}

This route asks whether one can choose a single isotropic frame of directions, depending on the
measure but not on the test function, so that normalized conditional line resampling has a
dimension-free spectral gap. It is a sufficient-condition route rather than a reformulation of
KLS. Its normalization is designed so that linear functions have exactly the Euclidean energy,
while the sharp one-dimensional log-concave Poincar\'e inequality compares the form from above
with the usual gradient form.

Let $\mu$ be a full-dimensional probability on $\R^d$ with finite second moment. For
$\theta\in S^{d-1}$, disintegrate along the affine lines parallel to $\theta$:
\[
  \int h(x)\dd\mu(x)
  =\int_{\theta^\perp}\int_\R h(z+t\theta)
    \dd\mu_{\theta,z}(t)\dd\bar\mu_\theta(z).
\]
Write $\sigma_{\theta,z}^2=\Var_{\mu_{\theta,z}}(T)$. An \emph{admissible frame} is an even
Borel probability $\rho$ on $S^{d-1}$ satisfying
\begin{equation}\label{eq:conditional-frame-isotropy}
  d\int_{S^{d-1}}\theta\theta^T\dd\rho(\theta)=I_d.
\end{equation}
For such a frame define
\begin{equation}\label{eq:conditional-fiber-form}
  \mathcal D_{\mu,\rho}(f)
  :=d\int_{S^{d-1}}\int_{\theta^\perp}
  \frac{\Var_{\mu_{\theta,z}}(f(z+T\theta))}
       {\Var_{\mu_{\theta,z}}(T)}
  \dd\bar\mu_\theta(z)\dd\rho(\theta).
\end{equation}
Null fibers and fibers with zero denominator contribute zero. The maximal form domain is
$\{f\in L^2(\mu):\mathcal D_{\mu,\rho}(f)<\infty\}$, using jointly measurable conditional
versions.

\begin{lemma}[Conditional-fiber form and factor-$4$ bridge, candidate]
\label{lem:conditional-fiber-form}
For a full-dimensional probability with density and finite second moment, and an admissible
frame, the quadratic form
\eqref{eq:conditional-fiber-form} is densely defined, closed, symmetric, Markovian, and
reversible. Its nonlocal self-adjoint generator is understood in form sense, and agrees with the
conditional pair-jump formula on the natural sufficient Bochner domain even when pointwise total
jump rates are infinite.

If $\mu$ is log-concave, then for every locally Lipschitz $f$ in the form domain,
\begin{equation}\label{eq:conditional-fiber-gradient-comparison}
  \mathcal D_{\mu,\rho}(f)
  \le4\int\abs{\nabla f}^2\dd\mu.
\end{equation}
Consequently, a form gap
$\Var_\mu(f)\le C\mathcal D_{\mu,\rho}(f)$ implies $C_P(\mu)\le4C$.
For every linear test $f_a(x)=a\cdot x$, one has
$\mathcal D_{\mu,\rho}(f_a)=\abs{a}^2$; hence linear tests have quotient one when $\mu$ is
isotropic. Every admissible frame has form gap one for the standard Gaussian, and the coordinate
frame has form gap one for a standardized product law.
\end{lemma}

\begin{question}[Conditional-fiber frame]
\label{q:conditional-fiber-frame}
For every full-dimensional isotropic log-concave $\mu$ on $\R^d$, does there exist an admissible
frame $\rho_\mu$, chosen independently of $f$, and a universal $C$ such that
\begin{equation}\label{eq:conditional-fiber-gap}
  \Var_\mu(f)\le C\mathcal D_{\mu,\rho_\mu}(f)
\end{equation}
for every $f$ in the maximal closed form domain?
\end{question}

An affirmative answer implies KLS by
\eqref{eq:conditional-fiber-gradient-comparison}. The order of quantifiers is essential:
$\rho_\mu$ may depend on $\mu$, but it must be fixed before the test $f$ is chosen.

\subsection{The simplex root frame: a rigorous failed implementation}

Let $P$ be uniform on the simplex $\Delta_{m-1}$, put $H_0=\mathbf1^\perp$ and
\[
  X=\sqrt{m(m+1)}\left(P-\frac1m\mathbf1\right)\in H_0.
\]
Then $X$ is isotropic in dimension $d=m-1$. The even $A_{m-1}$ root frame is
\[
  \theta_{ij}=\frac{e_i-e_j}{\sqrt2},
  \qquad
  \rho_{\rm root}=\frac1{m(m-1)}\sum_{i\ne j}\delta_{\theta_{ij}}.
\]
It satisfies
$d\int\theta\theta^T\dd\rho_{\rm root}=I_{H_0}$. If $\Var_{ij}$ denotes conditional
variance under redistribution of $(P_i,P_j)$ with their sum fixed, then
\begin{equation}\label{eq:conditional-root-form}
  \mathcal D_{\rm root}(f)
  =\frac{12}{m^2(m+1)}\sum_{i<j}
    \E\frac{\Var_{ij}(f)}{(P_i+P_j)^2}.
\end{equation}

\begin{proposition}[Vertex-cap obstruction for the root frame, candidate]
\label{prop:conditional-fiber-root-obstruction}
For $m\ge2$ and $\varepsilon\in(0,1)$, let $\eta_i=mP_i$ and
$A_{m,\varepsilon}=\{\eta_1>m-\varepsilon\}$. Then
\[
  p_{m,\varepsilon}:=\mu(A_{m,\varepsilon})
  =\left(\frac\varepsilon m\right)^{m-1},
\]
and the centered cap indicator belongs to the maximal root-form domain and satisfies
\begin{equation}\label{eq:conditional-root-cap}
  \frac{\mathcal D_{\rm root}(\one_{A_{m,\varepsilon}})}
       {\Var(\one_{A_{m,\varepsilon}})}
  \le
  \frac{12(m-1)}{(m+1)(m-\varepsilon)^2
  \left(1-(\varepsilon/m)^{m-1}\right)}.
\end{equation}
Consequently $\operatorname{gap}(\mathcal D_{\rm root})=O(m^{-2})$.
\end{proposition}

The vertex-cap failure is an $L^2$ phenomenon and does not descend to fixed polynomial degree.
The following candidate statement, isolated by the 2026-08-30 exact frame battery and staged in
a standalone dossier pending independent review (ledger status \emph{open}), pins the degree-two
root-frame pencil exactly.

\begin{lemma}[Root-frame degree-two pencil identity]
\label{lem:fiber-root-degree-two}
For every $m\ge3$, on the degree-$\le2$ quotient $V_{m,2}$ of $L^2$ of the isotropic uniform
simplex modulo constants, the $A_{m-1}$ root-frame pair form
$K(f,f)=\tfrac{12}{m^2(m+1)}\sum_{i<j}\E\bigl[\Var_{ij}(f)/(P_i+P_j)^2\bigr]$ and $G=\Var$
satisfy
\[
  \lambda_{\min}(K,G)=\frac{(m+2)(m+3)}{5m^2},
\]
attained exactly on the radial line $\R\,[\abs X^2-(m-1)]$. Consequently, granting the
certified identification of $K$ with the conditional-fiber root pencil, every degree-two dual
certificate in the all-frame min--max framework has objective at least $(m+2)(m+3)/(5m^2)>1/5$:
no vanishing degree-two refuter sequence exists, any fixed-degree polynomial dual refuter
requires degree at least three, and $\Lambda_{m,2}\ge(m+2)(m+3)/(5m^2)$. No upper bound on
$\Lambda_{m,2}$, no degree-three statement, and no claim about the full $L^2$ gap (where the
vertex-cap obstruction stands) is asserted.
\end{lemma}

The same failure mechanism has published prior art in negative-rate energy-exchange models.

\begin{proposition}[Negative-rate simplex exchange obstruction, imported]
\label{prop:sasada-negative-exchange}
For a symmetric-Dirichlet law of shape $a>0$, a complete-graph heat-bath form with pair rate
$(\eta_i+\eta_j)^s$ and $s<0$ has spectral gap at most $2^{-s}m^s$ in the exchange
normalization. In particular, at $a=1$ and $s=-2$, the corresponding $A_{m-1}$ root
conditional-fiber form has gap at most
\[
  \frac{48}{m(m+1)}.
\]
\end{proposition}

This is a published consequence of Sasada's vertex-cap argument
\cite{Sasada2015EnergyExchange}. Caputo proves a dimension-free gap for the unweighted flat
simplex heat bath, while Carlen--Posta--T\'oth prove uniform gaps for nonnegative exponents
$s\in[0,1]$ \cite{Caputo2008BinaryCollision,CarlenPostaToth2025Exchange}. The inverse-variance
root form lies at $s=-2$, outside those positive results.

The root calculation does \emph{not} refute Question~\ref{q:conditional-fiber-frame}. A general
permutation-invariant admissible frame may mix continuously many direction orbits, and no proof
shows that the root orbit is optimal. The decisive simplex alternative is an exact all-frame
dual certificate with objective tending to zero, or a uniform lower bound after optimizing over
all admissible frames. Floating finite sweeps and root-only tests decide neither statement.

\end{document}
